\documentclass[a4paper, landscape]{article}

% --- GEOMETRY & PAGE LAYOUT ---
% Set landscape mode with very small margins to maximize space
\usepackage[a4paper, landscape, top=0.5cm, bottom=0.5cm, left=0.75cm, right=0.75cm]{geometry}

% --- FONT & LANGUAGE (BABEL) ---
% Use fontspec for modern font handling (requires XeLaTeX/LuaLaTeX)
\usepackage{iftex}

% Use Unicode font features only for XeTeX or LuaTeX; fall back to pdfLaTeX safe fonts otherwise
\ifPDFTeX
  % Fallback for pdfLaTeX (no fontspec)
  \usepackage[T1]{fontenc}
  \usepackage[english]{babel}
  \usepackage{lmodern}
\else
  \usepackage{fontspec}
  \usepackage{xeCJK} % <-- added: Chinese support

  % Use babel for language settings
  \usepackage[english, bidi=basic, provide=*]{babel}
  \babelprovide[import, onchar=ids fonts]{english}

  % Prefer Noto Sans if available, otherwise fallback to Latin Modern or system fonts
  \IfFontExistsTF{Noto Sans}{\setsansfont{Noto Sans}}{%
    \IfFontExistsTF{Latin Modern Sans}{\setsansfont{Latin Modern Sans}}{\setsansfont{Arial}}%
  }

  \IfFontExistsTF{Noto Sans Mono}{\setmonofont{Noto Sans Mono}}{%
    \IfFontExistsTF{Latin Modern Mono}{\setmonofont{Latin Modern Mono}}{\setmonofont{Courier New}}%
  }

  % Chinese font
  \IfFontExistsTF{Noto Sans CJK SC}
    {\setCJKsansfont{Noto Sans CJK SC}}
    {%
      \IfFontExistsTF{FandolHei-Regular}
        {\setCJKsansfont{FandolHei-Regular}}
        {\setCJKsansfont{SimHei}}
    }

  % Use sans-serif as the default family for compact layout
  \renewcommand{\familydefault}{\sfdefault}
  \renewcommand{\CJKfamilydefault}{\CJKsfdefault}
\fi

% --- PACKAGES ---
\usepackage{multicol}      % For 5 columns
\usepackage{amsmath}       % For math
\usepackage{amssymb}       % For math symbols (like \rightarrow)
\usepackage{booktabs}      % For the new table
\usepackage{parskip}       % To control paragraph spacing
\setlength{\parskip}{1pt}  % Minimal spacing between paragraphs
\setlength{\parindent}{0pt} % No paragraph indentation
\usepackage[dvipsnames]{xcolor} % Added for coloring

% --- CUSTOM COMMANDS FOR COMPACT LAYOUT ---
% Command for main topics, adds a small vertical break
\newcommand{\topic}[1]{\par\medskip\textbf{#1}\par}

% Command for subtopics
\newcommand{\subtopic}[1]{\par\textbf{#1}}

% Command for code/keywords
\newcommand{\code}[1]{\texttt{#1}}

% Command for math
\newcommand{\mathsym}[1]{$#1$}

% --- APPLY STYLING ---
\definecolor{highlightcolor}{gray}{0.9} % Light gray for highlight

\renewcommand{\topic}[1]{%
  \par\smallskip\noindent
  \colorbox{highlightcolor}{%
    \parbox{\dimexpr\columnwidth-2\fboxsep\relax}{%
      \textbf{#1}%
    }%
  }%
  \par
}

\renewcommand{\subtopic}[1]{%
  \par\textbf{\color{blue}{#1}}%
}

% --- DOCUMENT ---
\begin{document}

% Remove all page numbering, headers, and footers
\pagestyle{empty}

% Set font to 5pt with 6pt line spacing (leading)
% This is extremely small, as requested.
\fontsize{4}{5}\selectfont

% Begin the 5-column layout
\begin{multicols}{5}

% --- COLUMN 1 ---

\topic{Week 0 Overview + Attention}
\par \textbf{ReLU vs. Sigmoid vs. SwiGLU：} \(\mathrm{ReLU}(x)=\max(0,x)\)，计算简单、正区间不易梯度消失，但负区间梯度为 0，可能出现 dead neuron；Sigmoid \(\sigma(x)=1/(1+e^{-x})\in(0,1)\)，适合概率/门控，但两端饱和、易梯度消失；SwiGLU 是平滑门控激活，约为 \(\mathrm{Swish}(xW_1)\odot(xW_2)\)，表达能力更强，但计算和参数更多。
\par \textbf{CNN：} 用同一个learned kernel/filter在不同空间位置滑动，利用local connectivity + weight sharing；每个 filter 学一种局部特征，深层逐渐组合成更复杂特征。
Convolution： kernel 在输入局部区域做加权求和；多个 kernel \(\Rightarrow\) 多个 output channels / feature maps。
Max Pooling： 对局部区域取最大值，降低空间分辨率和计算量，并提高对小幅平移的鲁棒性，但会丢失精确位置信息。
CNN vs. 全连接网络： CNN 利用空间局部性和参数共享，参数更少；全连接层中每个输入都连到每个神经元，对图像会产生大量参数。CNN可以绕开RNN的并行问题，同一层里所有位置可以并行计算。
ResNet / Skip Connection： \(y=F(x)+x\)；让网络学习 residual \(F(x)\)，同时给梯度提供直接传播路径，使深层网络更容易训练。
\par \textbf{RNN：} \(h_t=f(W_xx_t+W_hh_{t-1})\)；按时间顺序逐 token 处理，天然保留历史信息，但无法跨时间步并行，且容易出现梯度消失/爆炸和长距离依赖退化。
LSTM：在RNN中加入input/forget/output gates和cell state，帮助长期保存信息、缓解梯度消失，但仍然是串行计算。
序列CNN：可以并行处理不同位置并高效提取局部模式，但要建立长距离依赖通常需要堆很多层、扩大kernel或使用dilated convolution。
Transformer： 用 self-attention 代替 recurrence，使所有位置可以直接交互并可并行训练；RNN 需要 \(O(n)\) 个 sequential operations，而 self-attention 只需 \(O(1)\)，但完整 attention 的计算复杂度为 \(O(n^2d)\)。
为什么 Transformer 优于 RNN： RNN 必须逐 token 串行计算，而且长距离依赖较难学习；self-attention 可以直接连接远距离 token，最大路径长度更短，且训练更容易并行
\par \textbf{Loss:} L2 / MSE： \(L=\sum_i(y_i-\hat y_i)^2\)；对大误差惩罚更强，平滑、易优化，常用于 regression。
L1 / MAE： \(L=\sum_i|y_i-\hat y_i|\)；误差线性增长，因此比 L2 对 outlier 更鲁棒。Cross-Entropy： \(L=-\sum_k y_k\log p_k\)；one-hot 分类时 \(L=-\log p_{\text{true}}\)，即鼓励模型把高概率分配给正确类别；常用于分类和 next-token prediction。
L0 / L1 / L2 Regularization
- 正则化目标： \(L_{\text{total}}=L_{\text{data}}+\lambda R(w)\)；限制模型复杂度、减少 overfitting。
- L0： \(R(w)=\|w\|_0=\#\{w_i\neq0\}\)；直接鼓励稀疏模型，但离散、不可导，很难优化。
- L1： \(R(w)=\sum_i|w_i|\)；鼓励 sparsity，会把部分权重直接压到 0。
- L2： \(R(w)=\sum_iw_i^2\)；平滑地惩罚大权重，通常让权重趋近于 0 而不是直接变成 0，也常对应 weight decay。
- Dropout： 训练时随机把一部分 activation 设为 0，减少神经元之间的过度依赖，降低 overfitting；可近似看作训练很多不同子网络的 ensemble。
- Early Stopping： 监控 validation loss / performance，当验证集不再改善时停止训练，避免模型继续拟合训练集噪声，也是一种 regularization。
\par \textbf{SGD / Minibatch}
- Gradient Descent： \(\theta\leftarrow\theta-\eta\nabla_\theta L\)，其中 \(\eta\) 是 learning rate。
- SGD： 不用整个训练集计算精确梯度，而用单个样本或小部分样本估计梯度；单次更新便宜，而且梯度噪声有时有助于跳出不好的区域。
- Mini-batch SGD：\(
  g=\frac1B\sum_{i=1}^B\nabla_\theta L_i,
  \theta\leftarrow\theta-\eta g
  \)
  大 batch 梯度更稳定、GPU 并行效率高，但更耗显存；小 batch 噪声更大，但单次更新更便宜。
- Gradient Clipping： 限制梯度大小，防止 exploding gradients；常用 norm clipping 
- Adagrad： 给每个参数设置自适应 learning rate：\(
  G_{t,j}=\sum_{\tau=1}^{t}g_{\tau,j}^2,
  \theta_{t+1,j}
  =
  \theta_{t,j}
  -
  \frac{\eta}{\sqrt{G_{t,j}+\epsilon}}g_{t,j}
  \)
  经常更新的参数learning rate会变小，稀疏/少见特征相对获得更大步长；缺点是 \(G_t\) 一直累积，learning rate最后可能过小。
- RMSProp：用squared gradients的exponential moving average代替 Adagrad对全部历史梯度平方的累加，避免 Adagrad 的 learning rate 不断衰减到接近 0。Adam = Adagrad + RMSProp. AdamW ≠ Adam + ordinary L2 regularization；AdamW 的关键是 decouples weight decay from the gradient update
- BatchNorm： 在 mini-batch 内把 activation 归一化到近似均值 0、方差 1，再学习 scale 和 shift；能稳定训练、降低对初始化的敏感度，并常允许使用更大的 learning rate。BatchNorm 后的 \(\gamma,\beta\) 让网络重新学习合适的 scale 和 shift，不会被强制限制为 mean=0, variance=1。
\par n is the sequence length, d is the representation dimension, k is the kernel size of convolutions. Self-Attention 每层复杂度 \(O(n^2d)\)、可完全并行 \(O(1)\)、任意位置间路径长度 \(O(1)\)；RNN 为 \(O(nd^2)\)、必须串行 \(O(n)\)、长距离依赖路径 \(O(n)\)；CNN 为 \(O(knd^2)\)、可并行 \(O(1)\)、路径长度约 \(O(\log_k n)\)；restricted self-attention 将复杂度降为 \(O(rnd)\)，但最大路径长度增至 \(O(n/r)\)。
\par Byte-level BPE：先把文本变为 bytes，再反复合并常见相邻片段。优点是几乎无 OOV；缺点是罕见语言/奇怪文本可能token 数多。BERT 使用 WordPiece
\par Seq2Seq 的问题： 传统 encoder–decoder RNN 往往要把整个 source sentence 压缩进一个固定 hidden representation，长句时形成 bottleneck；attention 的核心改进是：生成每个 output token 时，直接关注 source 中相关位置，而不是只依赖一个压缩向量
\par Auto-regressive： decoder 一次生成一个 token，预测 \(y_t\) 时依赖此前的 \(y_{<t}\)；因此 inference 必须把之前生成的 token 作为下一步输入。原始 Transformer decoder 也是 auto-regressive。Embedding：把离散 token ID 映射成 \(d_{\text{model}}\) 维向量；embedding table参数量为
\(
|\text{Vocab}|\times d_{\text{model}}
\)
\par \textbf{Scaled Dot-Product Attention：}
\(
\mathrm{Attention}(Q,K,V)=
\mathrm{softmax}\left(\frac{QK^\top}{\sqrt{d_k}}\right)V
\)
若 sequence length 为 \(n\)，则 \(Q,K\in\mathbb R^{n\times d_k}\), \(V\in\mathbb R^{n\times d_v}\)；\(QK^\top\in\mathbb R^{n\times n}\) 是 token-to-token attention scores，softmax 后乘 \(V\) 得 \(n\times d_v\)。除以 \(\sqrt{d_k}\) 防止 dot product 随维度增大而过大、softmax 饱和、梯度过小。    
Q / K / V 的直觉： \(Q\)=“我在找什么”，\(K\)=“我有什么信息可供匹配”，\(V\)=“匹配后真正取走的信息”；attention 用 \(QK^\top\) 决定从各个 \(V\) 中取多少。
\par \textbf{Encoder–Decoder Attention：}Decoder 输出提供 \(Q\)，Encoder 输出提供 \(K,V\)，使 decoder 在生成每个 token 时可以关注输入序列中的相关位置。Encoder： 每层 = Multi-Head Self-Attention + Position-wise FFN，每个 sublayer 外有 residual + norm, 对整个输入序列生成 contextualized representations。Encoder self-attention 中 \(Q,K,V\) 都来自 input. Decoder： masked self-attention + encoder-decoder attention + FFN，自回归地逐 token 生成输出；位置 \(i\) 只能依赖之前已经生成的 token。 Decoder masked self-attention 中 \(Q,K,V\) 都来自 output prefix
\par \textbf{Multi-Head Attention：}
\(
\mathrm{head}_i=\mathrm{Attention}(QW_i^Q,KW_i^K,VW_i^V)
\)
\(\mathrm{MultiHead}(Q,K,V)=\mathrm{Concat}(\mathrm{head}_1,\ldots,\mathrm{head}_h)W^O
\)
不同 head 将输入投影到不同 representation subspaces，因此可同时学习不同位置关系/特征。若 \(d_{\text{model}}=512,h=8\)，则
\(
d_k=d_v=d_{\text{model}}/h=64.
\)
\(W_i^Q,W_i^K\in\mathbb R^{512\times64}\), \(W_i^V\in\mathbb R^{512\times64}\)；每个 head 输出 \(n\times64\)，8 个 head concat 后为 \(n\times512\)，再乘 \(W^O\in\mathbb R^{512\times512}\)，最终仍为 \(n\times512\)。多个 head 能同时关注不同 representation subspaces / positions；单个 head 的 weighted averaging 容易把不同关系混在一起。
\par \textbf{Positional Encoding：}self-attention 本身没有顺序概念，因此将 positional encoding 与 token embedding 相加；二者都为 \(d_{\text{model}}\) 维。使用不同频率的 sine/cosine：
\(
PE_{(pos,2i)}=\sin(pos/10000^{2i/d_{\text{model}}}),
\quad
PE_{(pos,2i+1)}=\cos(pos/10000^{2i/d_{\text{model}}}).
\)
\par \textbf{Position-wise FFN：}对每个 token 独立地、用同一组参数 应用同一个 MLP：
\(
FFN(x)=\max(0,xW_1+b_1)W_2+b_2.
\)
原 Transformer 中 \(d_{\text{model}}=512\rightarrow d_{ff}=2048\rightarrow512\)；“position-wise”不是不同 position 用不同网络，而是 same FFN applied separately to every position。Transformer Block： attention + feed-forward network，并在子层外加 residual connection 和 normalization
\par \textbf{Residual Connection / Add：}
\(
x+\mathrm{Sublayer}(x)
\)
让每层只需学习输入上的 incremental change，并给信息和梯度提供 shortcut，减少深层网络的 gradient problems
\par \textbf{LayerNorm：}对单个 token hidden vector 的 feature dimensions 做 normalization，再进行 learned scale/shift；主要作用是稳定 activation 和 optimization。Post-Norm：论文使用
\(
\mathrm{LayerNorm}(x+\mathrm{Sublayer}(x))
\)
即先Sublayer + residual，再 LayerNorm；这就是 Post-Norm.
Pre-Norm：
\(
x+\mathrm{Sublayer}(\mathrm{LayerNorm}(x))
\)
即先 LayerNorm，再进入 attention/FFN，最后 residual add；现代大型 Transformer 中很常见，通常让深层网络 optimization 更稳定。Pre-Norm 的 residual path 更接近 identity path，梯度可以更直接跨很多层传播
\par Causal Masking： decoder self-attention 中将未来位置的 attention logits 设为 \(-\infty\)，softmax 后权重变成 0，从而保证位置 \(i\) 只能看到 \(\le i\) 的 token，维持 autoregressive property。Shifted Right / Decoder Training： encoder 和 decoder 一起 end-to-end 训练；训练时 decoder 输入不是自己刚预测出的结果，而是 ground-truth target 向右移一位，从而可以并行预测下一 token；推理时没有 ground truth，因此必须把 previously generated output 逐步喂回 decoder。BERT： Transformer encoder-only；使用 multi-head self-attention、positional encoding 和 masking，双向attention，训练时通过 masked-token prediction 学习 bidirectional context，即预测内部被 mask 的 token. RoBERTa： 与 BERT same Transformer architecture，使用 BPE 而不是 BERT 的 WordPiece，并在更多数据上训练。GPT 是 decoder-style / causal Transformer，只能 attend 到过去 token，并 autoregressively 做 next-token prediction。

\topic{Week 1 Long Context + FlashAttention}
\par Long context 常见于 multi-document QA、AI coding agent 阅读整个 repository、长 chain-of-thought 等场景。标准 self-attention核心瓶颈为 \(O(n^2)\)。方法主要包括 sparse/local attention（Longformer）、approximate attention（Performer）、position/context extension（RoPE/YaRN）以及优化 exact attention implementation（FlashAttention）。

\par \textbf{Longformer：}用 sparse attention 替代 full self-attention，将local sliding-window attention与少量 task-motivated global attention结合；若每个 token 只 attend 宽度为 \(w\) 的邻域，则 complexity 从 full attention 的 \(O(n^2)\) 降为
\(
O(nw),
\)
当 \(w\ll n\) 时随 sequence length 近似 linear。

\par \textbf{Local Sliding-Window Attention：}每个 token 只 attend 周围固定 window 内的 tokens，而不是整个 sequence；优点是把 \(O(n^2)\) 降为 \(O(nw)\)，缺点是单层中远距离 tokens 无法直接交换信息。堆疊l层后receptive field约为l * w

\par \textbf{Receptive Field with Depth：}stacking multiple layers of local attention lets a token indirectly reach farther，类似 CNN 中 receptive field 随 depth 增长；因此即使每层只有 local attention，多层之后信息仍能逐步跨越较长距离传播。

\par \textbf{Dilated Sliding-Window Attention：}在固定 window budget 下跳过部分 intermediate positions，只 attend 间隔位置，从而在不显著增加每个 token attention 数量的情况下让每层 reach farther、扩大 receptive field。间隔d取样，receptive field约l*d*w，在相同连接数下更远

\par \textbf{Global Attention：}给少量 task-relevant tokens 全局 attention，使这些 token 可以 attend 整个 sequence，且其他 tokens 也可与其交互；例如 classification的\texttt{[CLS]}，QA的question tokens。Local attention primarily builds contextual representations，而 global attention allows the model to build full-sequence representations for prediction。local attention 和 global attention 使用 separate \(Q/K/V\) projection matrices，因为 local context modeling 与 whole-document aggregation 执行不同功能。

\par Longformer = Sliding Window + optional Dilation + Global Attention

\par Longformer Pretraining / LED：Longformer 可从 RoBERTa checkpoint 继续 masked-language-model pretraining，再 finetune；Longformer-Encoder-Decoder（LED）把 efficient Longformer attention 用到 encoder-decoder architecture，以支持 long-document seq2seq tasks，如 summarization。Performer: approximate full attention. Performer仍近似所有 token-pair interactions，用 FAVOR+（Fast Attention Via positive Orthogonal Random features）对 softmax attention 做 randomized kernel approximation，从而实现 approximately linear-time/space attention，

\par \textbf{RoPE：}Rotary Position Embedding 根据 token position 对 query/key vector 做 position-dependent rotation：
\(
q_i'=R(i)q_i,\quad k_j'=R(j)k_j.
\)
于是
\(
(q_i')^\top k_j'
=
q_i^\top R(i)^\top R(j)k_j
=
q_i^\top R(j-i)k_j,
\)
使 attention interaction 自然依赖 relative position \(j-i\)，而不是单纯 absolute position。

\par \textbf{Lost in the Middle：}研究问题不是 ``模型最多能输入多少 tokens''，而是模型能否robustly use information at arbitrary positions in a long context；论文控制 context length 和 relevant information 的 position 来测量性能变化。 Multi-Document QA Setup：输入由一个 question 和 \(k\) 个 documents 组成，其中 exactly one document contains the answer，其余 \(k-1\) 个为 distractors；保持内容不变、只移动 answer-containing document 在 context 中的位置，若模型能 robustly use long context，则性能应不应因 relevant information 的位置显著变化。\textbf{U-Shaped Performance：}模型通常在 relevant information 位于 context 开头时表现较好（primacy bias），位于结尾时也较好（recency bias），而位于中间时显著下降，larger context window does not imply robust context utilization。\textbf{Key-Value Retrieval：}即使在只需从 JSON-style key-value pairs 中找到指定 key 对应 value 的简单 synthetic retrieval task 中，一些模型仍在 relevant pair 位于 context middle 时性能下降并形成 U-shaped curve，说明问题不只是复杂 reasoning 导致。\textbf{Training Context Length：}encoder-decoder models 在输入长度处于 training-time sequence length 内时对 relevant-information position 相对 robust，但输入超过训练时见过的长度后也会出现 U-shaped performance；因此 lost-in-the-middle 不只是 nominal context-window limit 的问题。

\par \textbf{Instruction Tuning：}base language models（without instruction fine-tuning）也会出现 U-shaped performance curve，因此 lost-in-the-middle 并不能简单归因于 instruction fine-tuning。

\par \textbf{Query-Aware Contextualization：}把 query 放在 documents/key-value pairs 前后，让 representation 更明确地受到 query 影响；这种方法在 synthetic key-value retrieval 上可实现 near-perfect performance，但对 multi-document QA 的总体趋势影响有限. More Context Is Not Always Better. Evaluating Long-Context Models：只报告 maximum context window 不足以证明模型能有效使用 long context；更可靠的 evaluation 应改变 relevant information 在 context 中的位置，并检查 best-case 与 worst-case performance 是否接近，即性能是否对 position robust。

\par \textbf{FlashAttention：}FlashAttention 不改变 attention 的数学定义，而是重新组织 exact attention 的计算：
\(
S=QK^\top,\qquad
P=\mathrm{softmax}(S),\qquad
O=PV.
\)
核心目标不是 approximate attention，而是通过 IO-aware implementation 减少 GPU memory traffic. 减少 FLOPs 不一定减少实际运行时间, 现代 GPU 上 attention 等操作常受 memory access 而非 arithmetic throughput 限制

\par \textbf{GPU Memory Hierarchy：}GPU 中 HBM（High Bandwidth Memory）容量较大、off-chip、较慢，用于保存 model weights、activations、KV cache 等；SRAM 容量很小、on-chip、带宽高得多。CPU DRAM非常大但很慢。这里IO-aware指显式考虑不同 memory levels 的成本，核心目标是minimize reads/writes between HBM and SRAM.

\par \textbf{Why Standard Attention Is Slow：}standard attention 会 materialize 完整
\(
QK^\top\in\mathbb{R}^{N\times N},
\)
将其写入 HBM，再读取进行 softmax，再读取 softmax result 与 \(V\) 相乘；long sequence 时巨大的 \(N\times N\) intermediate matrix 被反复读写，runtime bottleneck 很大程度来自 data movement。

\par \textbf{Tiling：}将 \(Q,K,V\) 切分成能放入 fast SRAM 的 blocks，每次只 load 一部分数据，在 SRAM 内完成 score \(\rightarrow\) softmax \(\rightarrow\) value multiplication，再把需要的 output 写回 HBM；因此
\(\text{the full }N\times N\text{ attention matrix is never materialized in HBM}.
\) Q row block, KV column blocks. 每次只计算Sij=QiKjT，不需要写回HBM 扔掉

\par \textbf{Online Softmax：}
\(
\mathrm{softmax}(x_i)=\frac{e^{x_i}}{\sum_j e^{x_j}}
\)
看似必须先看到整个 row；FlashAttention 使用 online/streaming softmax，在逐 block 处理时维护 running maximum \(m\) 和 normalization sum \(\ell\)，并在新 block 改变最大值时 rescale 之前的部分结果，因此可以 block-by-block 计算但仍得到 exact softmax。
Numerically Stable Softmax：softmax 通常先减去 row maximum：
\(
m(x)=\max_i x_i,
f(x)_i=e^{x_i-m(x)},
\ell(x)=\sum_i f(x)_i,
\)
因此
\(
\mathrm{softmax}(x)=\frac{f(x)}{\ell(x)}.
\)
Online softmax 的关键是维护并更新这些 running statistics，而不保存完整 score row。

\par \textbf{Kernel：}GPU kernel 是在 GPU 上执行某个 operation 或一组 operations 的程序。Kernel Fusion：standard implementation 可能分别执行
\(
\text{matmul}\rightarrow\text{mask}\rightarrow\text{softmax}\rightarrow\text{dropout}\rightarrow\text{matmul},
\)
各 kernel 之间经常需要把 intermediate results 写回再从 HBM 读取；FlashAttention 将这些 attention operations fuse into one GPU kernel，使 intermediate data 尽量保留在 SRAM/registers 中，减少 HBM traffic。

\par \textbf{Backward Recomputation：}training backward 通常需要 forward intermediates；standard attention 会保存大的 \(N\times N\) attention matrix，而 FlashAttention 不保存它，只保存较小的 softmax normalization statistics，backward 时在 SRAM 中重新计算 attention blocks。虽然 recomputation 增加 FLOPs，但减少昂贵的 HBM reads，因此
\(
\text{more FLOPs}+\text{much less IO}
\Rightarrow
\text{lower wall-clock time}.
\)

FlashAttention extra O(n) memory

\par \textbf{FlashAttention IO Complexity：}若 head dimension 为 \(d\)、SRAM size 为 \(M\)，FlashAttention HBM-access complexity 为
\(
O\!\left(\frac{N^2d^2}{M}\right),
\)
而 standard attention 至少需要
\(
\Omega(Nd+N^2)
\)
HBM accesses

若主要时间花在 arithmetic operations 上，则为compute-bound，若主要时间花在 HBM reads/writes 上，则为memory-bound。
\(
\text{Arithmetic Intensity}
=
\frac{\text{number of arithmetic operations}}{\text{bytes of memory accessed}}.
\)

\par Training 时通常一次并行处理整个 sequence，因此 attention bottleneck 主要来自构造和保存 large \(n\times n\) attention intermediates；\textbf{KV Cache：}autoregressive decoding 中，每生成一个新 token，都需要让其 query attend 所有 previous tokens；为避免每一步重新计算历史 tokens 的 keys/values，模型缓存每层、每个 head、每个 previous token 的 \(K,V\)，称为 KV cache。其大小随 sequence length、batch size、number of layers 和 number of KV heads 线性增长。

\par \textbf{KV Cache Scaling：}若每层有 \(H\) 个 KV heads、每个 head dimension 为 \(d_h\)、sequence length 为 \(n\)、层数为 \(L\)，忽略 batch 和 datatype 时 KV cache roughly scales as
\(
2H d_h nL,
\)
其中 factor \(2\) 来自同时存储 keys 和 values；inference 通常分为 \textbf{prefill} 和 \textbf{decode}：prefill 对完整 prompt 并行计算 hidden states 和 KV cache，更像 training forward pass，计算量大且 matrix operations 较容易利用 GPU；decode 每次只处理一个新 token，需要不断读取 growing KV cache，因此 arithmetic intensity 较低、通常更加 memory-bandwidth-bound。

\par \textbf{MHA (Multi-Head Attention) KV Cache：}标准 MHA 中每个 query head 都有自己独立的 \(K,V\) projections，因此 \(H\) 个 heads 需要保存 \(H\) 组 key/value vectors：
\(
\text{KV cache}\propto 2H d_h nL.
\)
优点是每个 head 可以学习完全不同的 key/value representation、表达能力最强；缺点是 KV cache 最大、decoding memory bandwidth cost 很高。
\textbf{MQA (Multi-Query Attention)：}保留 \(H\) 个不同的 query heads，但所有 query heads share one K/V pair；因此 KV cache 相比 MHA 缩小约 \(H\times\)：
\(
H\text{ query heads}, 1\text{ KV head}.
\)
优点是显著降低 memory/storage 和 decoding cost；缺点是所有 heads attend over identical keys/values，head-specific flexibility 降低，可能造成 noticeable quality loss。
\textbf{GQA (Grouped-Query Attention)：}将 \(H\) 个 query heads 分成 \(G\) 组，每组共享一组 \(K,V\)：
\(
H\text{ query heads}, G\text{ KV heads}, 1<G<H.
\)
它是 MHA 与 MQA 的 compromise：相比 MHA，KV cache 约缩小 \(H/G\) 倍；相比 MQA，不同 groups 仍可学习不同 keys/values，因此通常保留更多 quality。
\textbf{MLA (Multi-Head Latent Attention)：}MLA 选择 compress, don't share：每个 token 的 K/V information 先压缩成一个 low-rank latent vector，只把 latent 存入 cache；attention 时再将 latent up-project/reconstruct 成 head-specific information，因此不同 heads 仍可从共享 compressed representation 中提取不同内容。
因此 cache 可比 typical GQA 更小，同时尽量恢复接近 MHA 的 expressiveness；代价是额外的 down-projection / up-projection compute，即典型的
\(
\text{spend extra compute to save memory}.
\)
Sparse Indexer for Decoding：应该 attend 哪些 previous tokens
Tiered KV-Cache Storage：除了缩小 KV cache，还可以把 KV cache 放到更便宜、更大的 memory hierarchy

\topic{Week 2 Computer Vision + Segmentation}
\par \textbf{Object Detection：}不仅判断图像中有什么 object，还需要定位 object 在哪里；典型输出为 bounding box + class + confidence。YOLO 的核心目标是 real-time, single-pass object detection，即一次 forward pass 同时完成 classification 和 localization。
\textbf{YOLO 核心思想：}传统 detection 方法通常先产生 candidate regions，再分别运行 classifier；YOLO instead ``frames object detection as a regression problem''，直接从整张 image pixels 回归 bounding boxes 和 class probabilities。整个 detection pipeline 是一个 single neural network，因此可以 end-to-end optimize。
\textbf{YOLO Grid：}将输入图像划分成 \(S\times S\) grid；如果某个 object 的 center 落在一个 grid cell 内，则该 cell 负责检测该 object。每个 cell 预测 \(B\) 个 bounding boxes，以及一组 \(C\) 个 conditional class probabilities \(P(Class_i\mid Object)\)。
每个 bounding box 预测
\(
(x,y,w,h,\mathrm{confidence}).
\)
其中 \(x,y\) 是 box center 相对于所属 grid cell 的位置，\(w,h\) 相对于整张图像；confidence 定义为
\(
P(Object)\cdot IoU_{\mathrm{pred}}^{\mathrm{truth}}.
\)
若 cell 中没有 object，则 confidence 应接近 0；若有 object，则希望 confidence 同时反映是否存在 object 和 localization quality。
整体 output tensor 为
\(
S\times S\times(B\cdot5+C).
\)
\textbf{Class-specific Confidence：}test time 将 conditional class probability 与 bounding-box confidence 相乘：
\(
P(Class_i\mid Object)P(Object)IoU
=
P(Class_i)IoU.
\)
因此最终 score 同时表示该 box 属于某一 class 的置信度以及 box localization 的准确程度。

\par \textbf{IoU / Intersection over Union：}
\(
IoU=\frac{|B_{\mathrm{pred}}\cap B_{\mathrm{gt}}|}
{|B_{\mathrm{pred}}\cup B_{\mathrm{gt}}|}.
\)
IoU 衡量 predicted bounding box 与 ground-truth box 的重叠程度；\(IoU=1\) 表示完全重合，\(IoU=0\) 表示完全不重叠。
\textbf{Responsible Bounding Box：}一个 cell 可以预测多个 boxes，但训练时只有与 ground-truth box 当前 IoU 最大的 predictor 对该 object ``responsible''，只有该 predictor 承担主要的 coordinate regression loss，从而让不同 box predictors 学习不同 box shapes/sizes。

\par \textbf{YOLO Loss：}YOLOv1 使用多部分 sum-squared error，包括 bounding-box coordinates、object confidence、no-object confidence 和 class probability：
\(
L=L_{\mathrm{coord}}+L_{\mathrm{obj}}+L_{\mathrm{noobj}}+L_{\mathrm{class}}.
\)
使用 \(\lambda_{\mathrm{coord}}=5\) 增强 localization error 的重要性，使用 \(\lambda_{\mathrm{noobj}}=0.5\) 降低大量 empty boxes 对 loss 的影响。
YOLO 不直接最小化 \(w,h\) 的平方误差，而比较
\(
(\sqrt w-\sqrt{\hat w})^2+
(\sqrt h-\sqrt{\hat h})^2.
\)
这样相同的 absolute error 对 small boxes 的影响比 large boxes 更大。
YOLO 在 training 和 inference 时都看到完整图像，因此 prediction 可以利用 global context；相比只观察 local region 的 region-proposal methods，更不容易把 background patch 错认成 object。
\textbf{YOLOv1 Architecture：}原始 YOLO backbone 主要由 convolution layers 提取 spatial features，最后用 fully connected layers 产生固定大小 detection tensor。backbone 是网络中负责从原始 image 提取 general visual features 的主体部分；后续 detection head 再利用这些 feature maps 预测 classes 和 bounding boxes。现代 YOLO architecture 通常可以理解为
\(
\mathrm{Backbone}\rightarrow\mathrm{Neck}\rightarrow\mathrm{Detection\ Head},
\)
其中 backbone 提取 features，neck 融合不同 scales 的 features，head 输出 detection predictions。先以224*224在ImageNet与训练convolutional layers，再升分辨率做detection fine-tuning。\textbf{优点：}single unified network 使 YOLO 简单、fast、end-to-end，并能 global reasoning。\textbf{局限：}有localization errors，难以精确定位 small objects；grid formulation 也会限制同一 cell 中多个靠得很近的 objects 的检测能力。

\par \textbf{Non-Max Suppression / NMS：}模型可能为同一个 object 产生多个重叠 bounding boxes；NMS 先保留最高-confidence box，然后删除与它 \(IoU\) 超过 threshold 的其他 boxes，并重复此过程。NMS 是 hand-designed greedy post-processing，而不是 learned component；相邻或高度重叠的真实 objects 有时也可能被错误 suppress。

\par \textbf{Anchor Boxes：}预设位置/尺宽/长宽比的参考框，模型预测相对偏移，覆盖形状多，但带来大量重复predictions。Backbone是提取通用features的主体网络。 \textbf{Focal Loss：}用于缓解 foreground/background 或 easy/hard examples 的严重不平衡，典型形式为
\(
FL(p_t)=-(1-p_t)^\gamma\log p_t.
\)
当 example 已经很容易、\(p_t\) 很高时，\((1-p_t)^\gamma\) 会降低它的 loss contribution，使模型更多关注 hard examples。代价gradient可能变小，训练收敛可能变慢。semantic segmentation 对每个 pixel 进行 classification，输出与图像空间位置对应的 segmentation map。Classic CNN连续pooling丢精确位置，sliding window patch classifier极慢

\par \textbf{U-Net 核心思想：}U-Net consists of a ``contracting path to capture context and a symmetric expanding path that enables precise localization''；左侧 encoder 逐渐提取更抽象的 context，右侧 decoder 恢复 spatial resolution，并通过 skip connections 恢复精细 localization information。
\textbf{U-Net Contracting Path / Encoder：}每一级重复
\(
3\times3\ \mathrm{Conv}+\mathrm{ReLU}
\rightarrow
3\times3\ \mathrm{Conv}+\mathrm{ReLU}
\rightarrow
2\times2\ \mathrm{MaxPool}.
\)
\(2\times2\) max pooling with stride \(2\) 使 \(H,W\) 减半；每次 downsampling 后 feature channels 加倍。
\textbf{Downsampling：}降低 spatial resolution，使后续 neurons 获得更大的 receptive field、能够利用更广的 context；代价是 precise spatial/localization information 会逐渐损失，因此 segmentation 不能只依赖 encoder。
\textbf{Feature Channel：}一个 feature map channel 表示某种 learned feature detector 的响应；增加 channels 并不是复制原 feature map，而是使用更多 learned filters 学习不同类型的 features。U-Net downsampling 时通常 spatial size 减半而 channel 数翻倍。Bottleneck layers压缩feature dimentionality让后续计算更便宜，同时帮助融合higher-level和lower-level information。

\par \textbf{U-Net Expansive Path / Decoder：}每一级先 upsample feature map，再使用 \(2\times2\) up-convolution 将 channel 数减半；随后与 contracting path 对应层的 high-resolution feature maps concatenate，再进行两次 \(3\times3\) convolution + ReLU。
\textbf{Skip Connection：}U-Net 将 encoder 中较早、spatial resolution 较高的 feature maps 直接 copy/concatenate 到对应 decoder layer。这样 decoder 同时获得 high-level semantic/context information 和 low-level precise spatial information，从而兼顾 recognition 与 localization。不是直接相加，而是concatenate+convolution.

\par \textbf{U-Net Output：}U-Net 没有 fully connected layers；最后使用 \(1\times1\) convolution 将每个 pixel 的 feature vector 映射到 class logits，再得到 pixel-wise class prediction，因此属于 fully convolutional network。

\par \textbf{U-Net Weighted Loss：}touching cells 之间只有很窄的 background region，普通 pixel-wise loss 很容易忽略这些 pixels；U-Net 因此对 object boundaries 赋予更大的 loss weight：
\(
w(x)=w_c(x)+w_0
\exp\left(
-\frac{(d_1(x)+d_2(x))^2}{2\sigma^2}
\right).
\)
其中 \(d_1,d_2\) 是该 pixel 到最近和第二近 object boundary 的距离，使两个 touching objects 之间的狭窄区域得到特别大的训练权重。

\par \textbf{SAM / Segment Anything：}SAM 将 segmentation 定义为 promptable segmentation：给定 image 和 segmentation prompt，模型输出一个 valid segmentation mask；目标是通过 prompting zero-shot transfer 到新的 image distributions 和 segmentation tasks。SAM Prompt：主要包括 point、bounding box 和 mask，论文也探索 text prompt。prompt 可以有歧义，因此模型必须能够处理 ambiguous prompts。

\par \textbf{SAM Architecture：} Image Encoder + Prompt Encoder -> Lightweight Mask Decoder -> Mask.
昂贵的 image encoder 对每张 image 只运行一次；同一个 image embedding 可以被多个不同 prompts 重复使用，因此 interactive segmentation 可以非常快。
\textbf{SAM Image Encoder：}使用 MAE-pretrained Vision Transformer (ViT) 处理高分辨率 image，将 image 转换成 dense image embedding；MAE pretraining 通过遮住大量 image patches、让模型重建 missing content 来学习 visual representations。
\textbf{SAM Sparse Prompt Encoder：}points 和 boxes 首先通过 positional encoding 表示其空间位置，再加上表示 prompt type 的 learned embedding，例如 foreground point 和 background point 使用不同 learned embeddings：
\(
e_{\mathrm{prompt}}
=
e_{\mathrm{position}}+e_{\mathrm{type}}.
\)
position 告诉模型 ``在哪里''，learned embedding 告诉模型 ``这个位置是什么意思''。
\textbf{SAM Dense Prompt Encoder：}mask 属于 dense prompt，因为它对大量 spatial locations 都提供信息；mask 通过 convolution 得到与 image embedding spatially aligned 的 embedding，再与 image embedding element-wise add。
\textbf{SAM Mask Decoder：}mask decoder 输入 image embedding、prompt embeddings 和 output token；通过 prompt self-attention 以及 prompt-to-image / image-to-prompt 双向 cross-attention 融合两类信息，之后 upsample image embedding，并用由 output token 产生的 dynamic linear classifier 计算每个 image location 的 foreground probability。
\textbf{Foreground / Background：}foreground 表示 prompt 希望 segment 的 target object pixels，background 表示不属于该 target 的 pixels。
\textbf{SAM Ambiguity-aware Prediction：}由于一个 prompt 可能存在多个合理 interpretation，SAM 对 ambiguous prompt 可以预测 multiple masks，而不是强迫模型输出所有合理答案的 average；训练时只对与 ground truth 最匹配的 predicted mask 计算主要 mask loss。

\par \textbf{SAM Mask Loss：}SAM 使用 focal loss 和 Dice loss 的组合：
\(
L_{\mathrm{mask}}
=
L_{\mathrm{focal}}+L_{\mathrm{dice}}.
\)
focal loss 重点学习 difficult foreground/background pixels；Dice loss 直接优化 predicted region 与 ground-truth region 的 overlap。
\textbf{Dice Score / Dice Loss：}
\(
Dice(P,G)=
\frac{2|P\cap G|}{|P|+|G|},
L_{\mathrm{Dice}}=1-Dice.
\)
Dice 更关心两个 segmentation regions 整体是否重合，因此在 foreground pixels 相对 background 很少时，比单纯 pixel accuracy 更合适。

\par \textbf{SAM Data Engine：}训练数据通过三阶段生成：assisted-manual stage 中 annotators 用 foreground/background clicks 配合 SAM 标注 masks；semi-automatic stage 中模型先生成 confident masks，人工补充遗漏 objects；fully automatic stage 中模型自动产生大量 masks。

\par \textbf{SAM 3 / Promptable Concept Segmentation：} Promptable Concept Segmentation (PCS) 给定 image/video 和 concept prompt，detect、segment，并在 video 中 track 所有 matching instances。concept 可以由 short noun phrase、image exemplar 或二者共同指定，text prompt 被限制为简单 noun phrases。SAM 的 point/box/mask prompt 通常告诉模型一个具体 object 在哪里；SAM 3 的 concept prompt 是 global 的，要同时完成 recognition + localization + segmentation。
\textbf{SAM 3 Architecture：}SAM 3 由共享 vision encoder 的 image-level detector 和 memory-based video tracker 构成。detector 负责在 image 中找到所有 matching concepts；tracker 负责在 video frames 中传播 masks 并保持 object identity。
SAM 3 Detector：detector 基于 DETR-style architecture，并 conditioned on text、geometry 和 image exemplars；它需要解决 open-vocabulary detection，即识别训练 label set 之外、由 prompt 动态指定的 visual concepts。
SAM 3 Presence Head：SAM 3 将 recognition 和 localization 解耦：独立 presence head 判断
\(
\text{``Is this concept present?''}
\)
而 detector localization 部分处理
\(
\text{``Where are its instances?''}
\)
这种 decoupling 对 difficult negative concepts 尤其有效。
Hard negative 是与 target concept 很相似、但实际上不满足 prompt 的 example；加入 hard negatives 可迫使模型学习更精确的 semantic decision boundary，而不是只依靠粗略 visual similarity。
\textbf{SAM 3 Video Tracking：}detector 应当 identity-agnostic，即只寻找所有 matching objects；tracker 则必须 identity-aware，区分不同 instances 并保持它们跨 frames 的 identity。
\textbf{SAM 3 Data Engine：}采用 human-and-model-in-the-loop pipeline；MLLM 可以作为 AI annotator 生成 noun phrases 和 hard negatives，也可以作为 AI verifier 检查 mask quality/exhaustivity；human annotators 主要集中修复困难错误，从而提高 annotation throughput。

\par \textbf{Spatial vs. Channel Attention：}spatial attention 在 location dimension 上选择 ``where''，即同一位置的 channels 一起增强/抑制；channel attention 在 feature dimension 上选择 ``what''，即同一个 channel 在所有 locations 一起增强/抑制。二者可以 sequential 或 parallel 使用。
\textbf{Parallel Spatial Attention：}``parallel'' 的核心意思是 attention branch 与部分普通 convolution/feature-processing branch 并行处理 features，之后再融合，而不是让所有 features 必须串行经过 attention；这样既保留原始/local convolution information，又用 spatial attention 强调 informative locations。
\par \textbf{Transposed Convolution / Up-convolution：}用于把 low-resolution feature map learnably 放大 spatial resolution；可以把普通 convolution 看成将 input positions 聚合到 output，而 transposed convolution 使用 convolution 的 transpose connectivity，把一个 input activation ``spread'' 到多个 output locations。是有 learnable kernel parameters 的 upsampling operation。increasing the stride increases both the height and width of the output. decoder 不只是 ``放大图像''，而是在恢复 resolution 的同时学习如何组合 semantic 和 localization information。

\par \textbf{TransUNet：}TransUNet ``replaces part of U-Net's CNN encoder with a transformer, keeps the skip-connection decoder''；Transformer encoder 用 self-attention 捕获 CNN 较难直接建模的 long-range/global context，而 U-Net-style decoder 和 skip connections 继续恢复精细 spatial detail。
Transformer 将 image 划成 patches/tokens 后可以获得 global relationships，但 patch-level representation 会损失部分 fine spatial detail；skip connections 把 encoder 较高 resolution 的 local features 直接传给 decoder，因此可以同时利用global context from Transformer + fine localization from U-Net. (TransUNet用Transformer捕获global context，同时保留UNet skip connections，把encoder的高分辨率空间细节传给decoder)

\par \textbf{SAM 2 Masklet：}SAM 2 的任务从 single-image mask 扩展为 video 中的 spatiotemporal mask。用户可以在任意 video frame 上提供 points、boxes 或 masks，定义一个 segment of interest，模型需要预测该 object 跨多个 frames 的 mask sequence，称为
\(
\text{masklet}.
\)
预测后还可以在后续 frames 加新 prompts 进行 iterative refinement。
mask 描述一个 object 在 单个 frame 中哪些 pixels 属于它；masklet 描述同一个 object 在 整个时间序列 中的 masks，因此同时包含 spatial segmentation 和 temporal consistency / tracking。
\textbf{SAM 2 Memory Attention：}Memory Attention 用过去 frames 的 features、past predictions 以及新的 prompts 来 condition 当前 frame features。SAM 2 stack \(L\) 个 transformer blocks；每个 block 执行
\(
\mathrm{SelfAttention}
\rightarrow
\mathrm{CrossAttention\ to\ Memory}
\rightarrow
\mathrm{MLP}.
\)
Self-attention 建模当前 frame 内部关系，cross-attention 从 memory bank 中读取过去 video information。
\textbf{SAM 2 Memory Encoder：}Memory encoder 将当前 predicted mask 先通过 convolutional module downsample，再与 image encoder 得到的 unconditioned frame embedding 做 element-wise sum，随后经过 lightweight convolution layers 融合：
\(
M_t
=
f_{\mathrm{conv}}
\left(
E_t+
\mathrm{Downsample}(\hat m_t)
\right).
\)
得到的 \(M_t\) 成为后续 frames 可以 cross-attend 的 memory representation。
\textbf{SAM 2 Memory Bank：}memory bank 保存 target object 的 past predictions，并维护 FIFO queues：最多 \(N\) 个 recent-frame memories，以及最多 \(M\) 个 prompted-frame memories；两类 memories 都以 spatial feature maps 存储。recent frames 提供最近 temporal context，prompted frames 保留用户明确纠正过的高价值 information。
ImageEncoder->MemoryAttention->MaskDecoder->MemoryEncoder->MemoryBank.
当前 frame 读取旧 memory 得到 mask，新 mask 再被编码进入 memory，形成 recurrent-like temporal processing loop。
如果为每个 object 独立维护 memory / inference state，则处理 \(K\) 个 objects 时 cost 会近似随 \(K\) 线性增长

\par \textbf{Scaling：}核心问题不是单纯 ``模型越大越好''，而是在 fixed compute budget 下决定 compute 应分配给 model size/ traning-date quantity/ training-data quality
因此 scaling 同时涉及 parameters、data、compute 和 data curation。
\textbf{Scaling Law：}经验上 model performance / loss 与 model size、dataset size 和 compute budget 之间呈近似 power-law relationship；常写为
\(
L(C)\approx L_\infty + aC^{-\alpha},
\)
因此在 log-log 坐标中，loss 与 compute 近似线性：
\(
\log(L-L_\infty)
=
\log a-\alpha\log C.
\)
这意味着增加 compute 通常带来 predictable but diminishing improvements，而不是突然饱和。随着 compute budget 增加，model size 和 training data 应近似按相同比例增加，而不是只把 parameters 做大. 当模型已经很大时，继续堆 parameters 未必比改善 training data 更有效。


% --- COLUMN 2 ---

\topic{Week 3 Diffusion Model + VLM}

\par \textbf{Visual Grounding：}建立 language concepts 与 visual objects/regions 之间的对应关系. ViLBERT 的核心思想是把 visual grounding 从某个特定 downstream task 中单独学习出来，变成一种可以 pretrain and transfer 的能力。

\par \textbf{CLIP：}Contrastive Language--Image Pre-training；不再训练模型预测固定的一组 object categories，而是从大量 \((image,text)\) pairs 中学习 image--text semantic alignment，使 natural language 本身成为 supervision，从而支持 zero-shot transfer。
\textbf{CLIP Dual Encoder：}CLIP 有独立的 image encoder \(f_I\) 和 text encoder \(f_T\)：
\(
I_i=\frac{f_I(x_i)}{\|f_I(x_i)\|},
T_j=\frac{f_T(t_j)}{\|f_T(t_j)\|}.
\)
两种 modality 分别编码，直到最终 embedding space 才通过 similarity 交互，因此属于 dual encoder / late fusion architecture。
\textbf{CLIP Similarity：}对 normalized image/text embeddings 使用 cosine similarity：
\(
s_{ij}=\frac{I_i^\top T_j}{\tau},
\)
其中 \(I_i^\top T_j\) 是 cosine similarity，\(\tau\) 是 learned temperature，用来控制 softmax distribution 的 sharpness。
\textbf{CLIP Contrastive Learning：}给定 batch 中 \(N\) 个正确的 image--text pairs，构造 \(N\times N\) similarity matrix；diagonal 的 \(N\) 个 pair 是 positives，其余 \(N^2-N\) 个组合是 negatives。目标是 increase similarity of correct pairs and decrease similarity of incorrect pairs，而不是直接生成 caption。
\textbf{CLIP Loss：}分别进行 image-to-text 和 text-to-image classification，再取平均：
\(
L=\frac12\left[
CE(s,\text{diag labels})
+
CE(s^\top,\text{diag labels})
\right].
\)
batch 越大，一个 example 可以比较的 negative examples 越多，因此 contrastive learning 通常受益于 large batch size。
\textbf{CLIP Image Encoder：}论文使用 modified ResNet 或 Vision Transformer (ViT)。modified ResNet 将最后的 global average pooling 替换为 attention pooling，使模型可以根据 learned query 对不同 spatial image features 加权聚合。
\textbf{CLIP Text Encoder：}使用 Transformer 编码 text；最高层 \([EOS]\) token 对应的 hidden state 被作为整段 text representation，然后通过 linear projection 和 normalization 映射到 shared embedding space。
\textbf{CLIP Shared Embedding Space：}image encoder 和 text encoder 将两种 modality 映射到同一个 vector space；语义相符的 image/text embeddings cosine similarity 应高，不相关 pair 的 similarity 应低。因此 shared space 中的 geometric distance 表达跨模态 semantic correspondence。

\par \textbf{CLIP Zero-shot Classification：}对每个 candidate class \(c\) 构造 natural-language prompt，例如 ``a photo of a \(c\)''，得到 text embedding \(T_c\)，然后把输入 image embedding \(I\) 与所有 \(T_c\) 比较：
\(
p(y=c\mid x)
=
\frac{\exp(s(I,T_c))}
{\sum_k\exp(s(I,T_k))}.
\)
不需要在 target dataset 上重新训练 classifier，因此称为 zero-shot transfer。传统 classifier 学固定 class weight vectors，而 CLIP 可以用 natural-language class descriptions 动态产生这些 vectors；因此 text encoder 可以理解为根据语言描述生成 zero-shot linear classifier weights。
\textbf{CLIP Prompt Engineering：}直接输入 class name（如 ``dog''）可能与 CLIP pretraining 时看到的 natural captions distribution 不一致，因此使用 ``a photo of a dog'' 等完整 prompt 通常效果更好；也可以 ensemble 多个 prompt templates。
\textbf{CLIP 的局限：}CLIP 的 image/text encoders 在中间层不发生 cross-modal interaction，只比较最终 global embeddings，因此适合 retrieval / zero-shot classification，但对于需要细粒度 token--region reasoning 的任务不如 cross-attention architecture 自然。

\par \textbf{ViLBERT：}Vision-and-Language BERT；一种 two-stream multimodal Transformer，visual stream 和 linguistic stream 分别处理 image 与 text，但在中间通过 co-attentional Transformer layers 交换信息，从而学习 task-agnostic joint visiolinguistic representations。
\textbf{Why Two Streams：}vision 与 language 的结构和 processing needs 不同，因此 ViLBERT 不从一开始把二者混在同一个 sequence 中；两种 modality 可以具有不同 network depth，并只在 selected layers 通过 co-attention 进行 sparse cross-modal interaction。
\textbf{ViLBERT Text Representation：}text side 基于 BERT；一个 token 的 input representation 是
\(
e_i=e_i^{token}+e_i^{position}+e_i^{segment},
\)
并使用 \([CLS]\)、\([SEP]\)、\([MASK]\) 等 special tokens。
\textbf{ViLBERT Visual Representation：}利用 pretrained object detector 提取 salient image regions；每个 region 由 visual feature 与 spatial/location information 表示，然后把这些 region representations 当作 visual tokens 输入 Transformer。

\par \textbf{ViLBERT Co-Attention：}核心是 ``exchanging key-value pairs in multi-headed attention''：visual queries attend to linguistic keys/values，同时 linguistic queries attend to visual keys/values：
\(
H_V'
=
\mathrm{Attention}(Q_V,K_L,V_L),
H_L'
=
\mathrm{Attention}(Q_L,K_V,V_V).
\)
这样 language information 可以改变 visual representations，visual information 也可以改变 linguistic representations。
\textbf{Self-Attention vs Co-Attention：}self-attention 负责 modality 内 reasoning，例如 words 之间或 image regions 之间的关系；co-attention 负责 cross-modal grounding。ViLBERT 同时保留两者，而不是所有层都做 co-attention。如果只有 co-attention，每一步 representation update 都依赖另一 modality；保留 ordinary self-attention layers 可以让每种 modality 先学习自身 internal structure，再在必要位置进行 cross-modal interaction。

\par \textbf{ViLBERT Masked Multimodal Modeling：}类似 BERT masked-language modeling，但同时应用于 language 和 vision；mask 一个 word 或 image region，再利用剩余 text + image context 预测其内容，使模型学习 intra-modal 与 cross-modal dependencies。
\textbf{ViLBERT Image--Text Alignment Task：}另一个 pretraining objective 是预测 image 与 text segment 是否 correspond；正例是真实 image-caption pair，负例是不匹配的 pair，从而训练 global multimodal alignment。
\textbf{ViLBERT Transfer：}pretrained ViLBERT backbone 只需 minor task-specific modifications 就可以 transfer 到 VQA、Visual Commonsense Reasoning、Referring Expressions 和 caption-based image retrieval 等任务。

\par \textbf{Early / Intermediate / Late Fusion：}Early fusion：较早把不同 modalities 合成统一 token sequence；Intermediate fusion：先分别处理，再在中间 layers 交互，ViLBERT 属于这一类；Late fusion：两种 modality 基本独立编码，到最后才结合，CLIP 属于这一类。

\par \textbf{Generative Model：}学习 observed data 的 probability distribution，使模型能够生成来自相同 distribution 的新 samples。\textbf{Autoencoder：}encoder 把高维 input 压缩成 latent representation：
\(
z=E(x),
\)
decoder 再重建：
\(
\hat x=D(z).
\)
训练时通过 reconstruction loss 使 \(D(E(x))\approx x\)；latent \(z\) 应保存重要 information，同时去掉 input 中的 redundancy。
\textbf{Latent Space：}encoder 学到的 compact representation space；不是原始 pixel space，而是用较低维向量表示 input 中最重要的 semantic / perceptual features，因此在 latent space 上进行 computation 通常比直接处理所有 pixels 更便宜。

\par \textbf{VAE / Variational Autoencoder：}普通 autoencoder 对每个 \(x\) 学一个 deterministic latent vector \(z\)，而 VAE ``learns a latent distribution, instead of a latent representation''：
\(
q_\phi(z\mid x)
=
\mathcal N(\mu(x),\sigma^2(x)I).
\)
encoder 输出 \(\mu,\sigma\)，再从这个 distribution sample \(z\) 后交给 decoder。
\textbf{VAE Reparameterization Trick：}
\(
z=\mu+\sigma\odot\epsilon,
\qquad
\epsilon\sim\mathcal N(0,I).
\)
把 random sampling 写成 deterministic function of \(\mu,\sigma,\epsilon\)，使 gradients 可以通过 \(\mu,\sigma\) backpropagate。
\textbf{VAE Loss：}
\(
L
=
L_{\mathrm{reconstruction}}
+
D_{KL}
\big(q_\phi(z\mid x)\,\|\,\mathcal N(0,I)\big).
\)
Reconstruction term 让 decoded image 接近 input；KL term 让 latent distributions 接近 standard Gaussian prior，使 latent space smooth、continuous 且容易 sampling。
\textbf{Why VAE Can Generate：}因为训练时要求
\(
q_\phi(z\mid x)\approx\mathcal N(0,I),
\)
所以 inference 时可以直接 sample
\(
z\sim\mathcal N(0,I)
\)
再通过 \(D(z)\) 生成新 image，而不一定需要对应某个 input image。

\par \textbf{Diffusion Model：}训练时逐渐向真实 data 加 Gaussian noise，并训练 neural network 学会反向 denoise；generation 时从 pure Gaussian noise 开始，通过多步 denoising 得到新 sample。可以理解为 ``a sequential application of denoising autoencoders''。

\par \textbf{Forward Diffusion / Noising Process：}逐步向 \(x_{t-1}\) 添加少量 Gaussian noise：
\(
q(x_t\mid x_{t-1})
=
\mathcal N
\left(
\sqrt{1-\beta_t}\,x_{t-1},
\beta_tI
\right),
\)
其中 \(\beta_t\) 控制第 \(t\) 步加入多少 noise；经过足够多 steps 后
\(
x_T\approx\mathcal N(0,I).
\)

\par \textbf{Diffusion Noise Schedule：}定义
\(
\alpha_t=1-\beta_t,
\bar\alpha_t=\prod_{s=1}^{t}\alpha_s,
\)
则不需要真的一步步执行 forward process，可以直接从 \(x_0\) sample 任意 timestep：
\(
x_t
=
\sqrt{\bar\alpha_t}x_0+
\sqrt{1-\bar\alpha_t}\epsilon,
\epsilon\sim\mathcal N(0,I).
\)

\par \textbf{Diffusion Training Objective：}随机 sample clean image \(x_0\)、timestep \(t\) 和 Gaussian noise \(\epsilon\)，构造 noisy \(x_t\)，让 neural network 预测加入的 noise：
\(
L_{\mathrm{simple}}
=
\mathbb E_{x_0,t,\epsilon}
\left[
\|\epsilon-\epsilon_\theta(x_t,t)\|_2^2
\right].
\)
因此模型通常不是直接预测最终 clean image，而是学习 ``given this noise level, what noise should be removed?''。

\par \textbf{Reverse Diffusion / Sampling：}generation 从
\(
x_T\sim\mathcal N(0,I)
\)
开始，然后反复执行 learned reverse denoising：
\(
x_T\rightarrow x_{T-1}\rightarrow\cdots\rightarrow x_1\rightarrow x_0.
\)
因为需要 sequential model evaluations，diffusion inference 通常比 single-pass generator 慢。

\par \textbf{Timestep Embedding：}同一个 denoising network 要处理从 nearly-clean 到 almost-pure-noise 的不同 \(x_t\)，因此把 timestep \(t\) 编码成 embedding 并输入网络，让模型知道当前 noise level。
\textbf{Why Diffusion Is Expensive：}pixel-space diffusion 每一步都需要在完整 high-dimensional RGB image 上运行 network；training 需要 repeated forward/backward evaluations，inference 也需要 tens 到 hundreds 甚至更多 sequential denoising steps，因此 computation 很大。

\par \textbf{Latent Diffusion Model (LDM)：}不是直接在 RGB pixels 上做 diffusion，而是先用 pretrained autoencoder：
\(
z_0=E(x),
\)
然后在 latent space \(z\) 中执行 diffusion；最终 denoise 得到 \(z_0\) 后：
\(
\hat x=D(z_0).
\)
这样 diffusion network 处理的是更小的 latent representation，而不是所有 pixels，因此显著降低 training / inference computation。
Pixel DM：
\(
x_t\in\mathbb R^{H\times W\times3}
\)
直接在 RGB space denoise；LDM：
\(
z_t\in\mathbb R^{h\times w\times c},
 h<H,\;w<W,
\)
在 compressed latent space denoise，再通过 decoder 恢复 pixel image。

\par \textbf{Perceptual Compression：}LDM 先用 autoencoder 去掉大量 perceptually irrelevant / high-frequency information，同时尽量保留 visually important details；这样 diffusion model 不需要花大量 compute 去建模几乎看不见的 pixel-level details。
\textbf{Semantic Compression：}在 autoencoder 完成 perceptual compression 后，diffusion model 在 latent space 中学习更 high-level 的 semantic / conceptual composition，例如 objects、layout、style 等。
\textbf{LDM Compression Trade-off：}latent 压缩太少则 compute reduction 不明显；压缩太强则丢失 reconstruction details。LDM 的设计目标是在 complexity reduction 与 detail preservation 之间找到 near-optimal trade-off。

\par \textbf{Stable Diffusion High-level Pipeline：}
\(
\mathrm{text}
\rightarrow
\mathrm{text\ encoder}
\rightarrow c,
\)
\(
x
\xrightarrow{E}
z_0
\xrightarrow{\mathrm{add\ noise}}
z_t
\xrightarrow[\mathrm{condition}\ c]{\mathrm{U\!-\!Net}}
\hat\epsilon,
\)
重复 denoising 得到 \(z_0\)，再
\(
D(z_0)\rightarrow\mathrm{RGB\ image}.
\)

\par \textbf{U-Net in Diffusion：}U-Net 接收 noisy latent/image \(z_t\)、timestep embedding 和 conditioning information；downsampling path 逐渐降低 spatial resolution、提取更 global features，upsampling path 恢复 spatial resolution，skip connections 将对应 resolution 的 earlier features 直接传给 decoder side，从而保留 local/fine-grained details。
\textbf{Text Conditioning：}text prompt 先由 text encoder 转换为 token embeddings \(c\)；这些 text embeddings 本身不会直接生成 pixels，而是作为 conditioning information，在每一步 denoising 中告诉 model 应朝什么 semantic direction 去除 noise。
\textbf{Cross-Attention in Latent Diffusion：}visual/latent representations 提供 queries，conditioning text embeddings 提供 keys 和 values. 因此不同 spatial locations 可以 attend 到 prompt 中不同 words，使 text information 直接控制 denoising。
conditioning 不需要与 image/latent 具有相同 spatial dimensions；只要先 encode 成 token representations，就可以作为 \(K,V\) 输入 cross-attention，因此同一 LDM architecture 可以接受 text、class labels、bounding boxes、layouts 等不同 conditioning inputs。

Diffusion model 最终产生的是 latent \(z_0\)，不是 RGB pixels；decoder \(D\) 将
\(
z_0\rightarrow x.
\)
若删除 decoder，仍然可以得到 latent representation，但无法直接得到最终可观看的 RGB image。
Autoencoder stage 是 reconstruction task：
\(
D(E(x))\approx x.
\)
但 diffusion stage 不是要求从 noise 最后还原某一张特定 training image；它学习 latent data distribution，generation 从 random noise 开始，只需要生成 plausible samples satisfying the condition，而不是与某个原始 \(x\) 相等。

\par \textbf{Guidance：}在 diffusion sampling 过程中改变 denoising direction，使生成结果更符合某个 condition；diffusion formulation 的优点之一是可以利用 guiding mechanism 控制 generation，而不需要为每一个新的 prompt 重新训练 generator。
\textbf{Classifier-Free Guidance (CFG)：}训练时有一部分 examples 故意去掉 condition，因此同一个 model 同时学会 unconditional prediction
\(
\epsilon_\theta(x_t,\varnothing)
\)
和 conditional prediction
\(
\epsilon_\theta(x_t,c).
\)
sampling 时组合二者：
\(
\hat\epsilon
=
\epsilon_\theta(x_t,\varnothing)
+
w\left[
\epsilon_\theta(x_t,c)
-
\epsilon_\theta(x_t,\varnothing)
\right].
\)

\par \textbf{CFG Intuition：}
\(
\epsilon_\theta(x_t,c)-\epsilon_\theta(x_t,\varnothing)
\)
近似表示 ``condition \(c\) 希望 denoising direction 改变多少''；guidance scale \(w\) 放大这个方向。较大的 \(w\) 通常提高 prompt adherence / sample fidelity，但过大会减少 diversity，甚至产生 artifacts。

\par \textbf{DALL-E 2 / unCLIP：}DALL-E 2 不直接从 text embedding 一步生成 image，而是显式先预测 CLIP image latent：
\(
\mathrm{text}
\rightarrow
\mathrm{CLIP\ text\ embedding}
\rightarrow
\mathrm{prior}
\rightarrow
\mathrm{CLIP\ image\ embedding}
\rightarrow
\mathrm{diffusion\ decoder}
\rightarrow
\mathrm{image}.
\)
论文把整个 architecture 称为 unCLIP，因为 generation 可以理解为 ``invert the CLIP image encoder''。

\par \textbf{DALL-E 2 Prior：}prior 学习
\(
P(z_i\mid y),
\)
其中 \(y\) 是 caption / text condition，\(z_i\) 是对应 image 的 CLIP image embedding；即先把 language semantics 转换成一个 plausible image representation，而不是直接预测 pixels。
CLIP text embedding 和 CLIP image embedding 虽然位于 shared semantic space、相互对齐，但并不是完全相同的 representations/distributions；prior 学习
\(
\mathrm{text\ representation}\rightarrow
\mathrm{possible\ image\ representation},
\)
再让 image decoder 处理它熟悉的 image embedding。
\textbf{DALL-E 2 Diffusion Prior：}论文比较 autoregressive prior 与 diffusion prior，发现 diffusion prior 更 computationally efficient，并得到更好的 sample quality；因此 production-style unCLIP 使用 diffusion prior。

\par \textbf{DALL-E 2 Decoder：}diffusion decoder 学习
\(
P(x\mid z_i,y),
\)
即根据 CLIP image embedding \(z_i\)（以及可选 text condition \(y\)）产生 actual image \(x\)。
\textbf{DALL-E 2 Probability Factorization：}整个 text-to-image model 可写为
\(
P(x\mid y)
=
P(x,z_i\mid y)
=
P(x\mid z_i,y)P(z_i\mid y).
\)
先由 prior sample image embedding \(z_i\)，再由 decoder sample actual image \(x\)。

\par \textbf{Why DALL-E 2 Is Stochastic：}同一句 caption 可以对应很多不同但正确的 images，因此 prior 不应该学习唯一 deterministic mapping，而是
\(
z_i\sim P(z_i\mid y).
\)
同样，一个 CLIP image embedding 并不保存所有 pixel-level details，因此 decoder 也可以从同一个 \(z_i\) 生成多个 semantically similar images。
\textbf{Image Variation in DALL-E 2：}给定 input image：
\(
x\rightarrow E_{\mathrm{CLIP}}(x)=z_i,
\)
然后重复 sample
\(
x'\sim P(x\mid z_i).
\)
因为 CLIP embedding 保存主要 semantics/style、但丢弃一些 non-essential details，因此不同 samples 能保持 overall meaning/style，同时改变 pose、background、minor details 等。

\par \textbf{Language-guided Image Manipulation：}CLIP 的 image 和 text 位于 joint embedding space，因此可以构造 semantic direction，例如
\(
d=
E_T(\text{``smiling''})
-
E_T(\text{``neutral''}),
\)
将 image embedding 沿 \(d\) 移动后再 decode，就可以进行 zero-shot language-guided semantic editing。

\par \textbf{DALL-E 2 Upsampling：}base diffusion decoder 先生成低分辨率 image，再使用 additional diffusion upsamplers：
\(
64\times64
\rightarrow
256\times256
\rightarrow
1024\times1024.
\)
upsamplers 主要负责增加 spatial resolution 和 fine detail，而不是重新决定 high-level semantics。


\topic{Week 4 Agentic AI}
\par \textbf{RAG (Retrieval-Augmented Generation)：}模型先从外部文档库 retrieve relevant documents，再把 retrieved context 连同 query 一起交给 generator 生成答案：
\(
x \rightarrow \mathrm{Retriever} \rightarrow z_1,\ldots,z_K \rightarrow \mathrm{Generator} \rightarrow y.
\)
核心思想是让模型访问训练参数之外、可更新的 external knowledge；RAG 可以减少仅靠 parametric memory 带来的 outdated/missing knowledge，但 不能保证消除 hallucination，因为 generator 仍可能误读、忽略或错误组合 retrieved information。

\par \textbf{DPR Retriever：}原始 RAG 使用 Dense Passage Retrieval (DPR)，分别用两个 BERT encoder 编码 query 和 document：
\(
q(x)=\mathrm{BERT}_q(x),\quad d(z)=\mathrm{BERT}_d(z),
\)
并用 inner product 衡量相关性：
\(
p_\eta(z\mid x)\propto \exp(d(z)^\top q(x)).
\)
DPR 是 bi-encoder，因此 document embeddings 可以提前计算并建立 index；query 到来后只需编码 query 再进行 retrieval。MIPS = Maximum Inner Product Search，用于高效找到 inner product 最大的 top-\(K\) documents。

\par \textbf{RAG Generator：}原始 RAG generator 使用 BART-large；给定 query \(x\) 和 retrieved document \(z\)，生成答案：
\(
p_\theta(y\mid x,z).
\)
Retriever 提供 external knowledge，generator 负责基于 query + document 做 conditional generation。

\par \textbf{RAG as Latent-Variable Model：}训练数据只有 query-answer pairs \((x,y)\)，并不需要标注“哪个 document 是正确证据”；retrieved document \(z\) 被视为 latent variable，对不同 document 的生成概率进行 marginalization：
\(
p(y\mid x)=\sum_z p_\eta(z\mid x)p_\theta(y\mid x,z).
\)
因此 retriever 和 generator 可以通过最终 answer likelihood 联合学习。\textbf{RAG-Sequence：}假设同一个 latent document \(z\) 支撑整条 output sequence：
\(
p(y\mid x)=
\sum_{z\in \mathrm{top}K}
p_\eta(z\mid x)
\prod_i p_\theta(y_i\mid x,z,y_{<i}).
\)
同一个 candidate document 在整条 sequence 内保持不变。
\textbf{RAG-Token：}允许每个 generated token 从不同 retrieved documents 获得支持：
\(
p(y\mid x)=
\prod_i
\sum_{z\in \mathrm{top}K}
p_\eta(z\mid x)
p_\theta(y_i\mid x,z,y_{<i}).
\)
比 RAG-Sequence 更 flexible，因为不同 token 可以依赖不同 evidence。

常见方案是
\(
\text{retrieve first}+\text{reason over retrieved text with long context}.
\)
RAG 的优势是减少 irrelevant context 和 context-window cost。

\par \textbf{Toolformer：}让 LM 通过 self-supervised training 学会 \textit{which APIs to call, when to call them, what arguments to pass, and how to incorporate the results}；只需要每个 API 少量 demonstrations，然后由模型自动生成大量 tool-use examples。
\textbf{Toolformer Training Pipeline：}
\(
\text{Sample API Calls}
\rightarrow
\text{Execute API Calls}
\rightarrow
\text{Filter Useful Calls}
\rightarrow
\text{Fine-tune LM}.
\)
模型先自己 propose potential calls，真正执行获得 tool results，再只保留能够提高后续 token prediction 的 calls，最后把它们插回 training corpus 进行普通 LM fine-tuning。

\par \textbf{Toolformer API Representation：}一个 API call 表示为
\(
c=(a_c,i_c),
\)
其中 \(a_c\) 是 API name，\(i_c\) 是 input；linearized call：
\(
e(c)=\langle API\rangle a_c(i_c)\langle/API\rangle,
\)
带 result：
\(
e(c,r)=\langle API\rangle a_c(i_c)\rightarrow r\langle/API\rangle.
\)
tool call 和普通 text 一样被插入 token sequence 中.

\par \textbf{Toolformer Call Sampling：}给定原始 text \(x=x_1,\ldots,x_n\)，模型在候选位置 \(i\) 判断是否开始生成 API call：
\(
p_i=p_M(\langle API\rangle\mid x_{1:i-1}).
\)
只在概率足够高的位置 sample candidate tool calls，因此“在哪里调用”和“调用什么”都由 LM 自己提出，而不是人工逐条标注。

\par \textbf{Toolformer Filtering：}判断一个 tool call 是否有用的标准是：加入 tool result 后，模型预测后续 tokens 的 loss 是否明显下降。定义
\(
L_i(z)=-\sum_{j=i}^{n}w_{j-i}\log p_M(x_j\mid z,x_{1:j-1}).
\)
保留 API call 当
\(
L_i^- - L_i^+ \ge \tau_f,
\)
其中 \(L_i^+\) 是提供 API result 后的 loss，\(L_i^-\) 是不用有效 result 的 baseline；即keep a tool call if its result helps predict future tokens.

\par \textbf{Toolformer Self-Supervision：}Toolformer 不需要人工标注“这个位置应该调用哪个 tool”；LM 先用少量 demonstrations 自动 annotate 大量文本，再用 future-token loss 判断 tool result 是否真正有用，因此 usefulness 本身成为 training signal。
\textbf{Toolformer Limitation：}Toolformer 仍需要少量 examples 定义 API usage format，且主要研究 relatively simple textual API calls；API input/output 需要能表示成 text，它没有完整解决现代 long-horizon agent 所需的 planning、retry、permissions、memory、verification 等系统问题。

\par \textbf{Agent：}与一次性的
\(
x\rightarrow \mathrm{LLM}\rightarrow y
\)
不同，agent 会在 environment 中反复采取 action 并读取 observation：
\(
s_t\rightarrow a_t\rightarrow o_t\rightarrow s_{t+1}.
\)
因此 agent 的输出会影响未来输入，是一个 feedback loop，而不是单次 inference。
\textbf{ReAct：}ReAct = Reason + Act，将 reasoning traces 和 task-specific actions interleave：
\(
\text{Thought}\rightarrow\text{Action}\rightarrow\text{Observation}\rightarrow\text{Thought}\rightarrow\cdots
\)
Reason 决定“下一步做什么、为什么”；Act 调用 tool 或执行 action；Observe 读取结果，然后基于新 information 再 reasoning。
\textbf{Why ReAct：}reasoning traces 可以帮助 agent induce, track, and update action plans、处理 exceptions；actions 则让模型从 external sources/environment 获得新的 information。reasoning guides actions; observations ground reasoning
而不是纯粹依赖内部 reasoning 或盲目执行 actions。

CoT 是 reason-only：可以进行 multi-step reasoning，但不能主动获得 external evidence，因此早期 hallucination 可能沿 reasoning chain 继续 propagate；ReAct 可以通过 search/tool observation 检查和修正 intermediate beliefs。Act-only 可以不断调用 tool/environment，但缺少显式 reasoning state，因此可能不知道下一步为什么要这样做、重复无效 actions 或 search wrong entity；ReAct 用 Thought 保存当前 subgoal、已知信息和下一步 plan。

\par \textbf{ReAct Failure Modes：}Reasoning trace 本身仍可能 hallucinate；agent 也可能进入 repetitive action loops，例如重复同一个 search/action 却不改变策略。常见 harness-level mitigation 包括 step limit、loop detection、replanning、validation 和 stopping rule。

\par \textbf{Harness：}agent harness 是围绕 LLM 的 execution/infrastructure layer，负责 context construction、tool interaction、orchestration、verification、retry、stopping 等：
\(
\text{Agent System}=\text{LLM}+\text{Harness}.
\)
同一个 model 使用不同 harness 时 performance 可能显著不同，因此 agent benchmark score 不是单纯的 model property。
\textbf{Harness as Closed-Loop Controller：}把 agent 看成 closed-loop dynamical system：state \(s_t\) 包含 current context + persistent memory，LLM 是 stochastic policy
\(
\pi_\theta(a_t\mid s_t),
\)
生成 action \(a_t\)；tool/environment 返回 observation \(o_t\)，harness 再构造下一步 state。论文因此把 harness 视为 controller，而 model 只是被 controller 包围的 policy。
\textbf{Harness Decisions：}每个 agent step 都包含 harness choices：task/system prompt 如何构造、tool catalogue 如何加入、tool call 如何 parse/execute、error 如何 retry、context 如何 compress、何时 stop/escalate。因此最终 benchmark score 是 model 与这些 harness choices 的 joint outcome。

\par \textbf{Binding Constraint Thesis：}Harness paper 主张，在 long-horizon tasks、且比较 capability 相近的 frontier models 时，performance variance 可能更多来自 harness configuration 而不是 model choice；因此不 disclosure harness 的 cross-model leaderboard 可能把 harness-level improvement 错误归因于 model。
\textbf{Locked-Harness Evaluation：}比较 models 时固定同一个 harness \(H\)：
\(
B(M_1,H),B(M_2,H),\ldots
\)
这样 model 是主要变化变量，可以更合理地 isolate model effect。

\par \textbf{Factorial Evaluation：}同时系统改变 model \(M_i\) 和 harness \(H_j\)，评估
\(
B(M_i,H_j),
\)
从而把 performance variation 分解为 model effect、harness effect 和 model-harness interaction，而不是把所有 gain 都归给 model。

\par \textbf{Harness Variance vs. Model Variance：}
\(
HV(M_i)=\mathrm{Var}_{H}[B(M_i,H)],
\)
衡量固定 model、改变 harness 时的 performance variance；
\(
MV(H_j)=\mathrm{Var}_{M}[B(M,H_j)],
\)
衡量固定 harness、改变 model 时的 variance。若
\(
HV\gg MV,
\)
说明当前 benchmark 更受 harness choice 影响。

\par \textbf{Ranking Reversal：}若
\(
B(M_1,H_1)>B(M_2,H_1)
\)
但
\(
B(M_1,H_2)<B(M_2,H_2),
\)
则 model ranking 会因 harness 改变而反转，因此单独说“model \(M_1\) 比 \(M_2\) 强”并不充分，必须说明使用的 execution harness。
\textbf{Context Management：}agent trajectory 很长时，harness 必须决定哪些旧信息保留、summarize、retrieve 或删除；不好的 context management 会导致 important constraints/files 从 effective context 中消失，即 context drift。
\textbf{Retry / Recovery：}tool call 可能 timeout、格式错误或返回 invalid result；harness 需要定义 retry、fallback、replanning、checkpoint/rollback 等 recovery policy，否则 agent 容易在一次局部错误后不断累积错误。
\textbf{Verification：}Harness 可以在 intermediate actions 或 final output 后运行 tests、schema validation、self-check 或 external verifier；verification 是 model 外部 reliability mechanism，不应该被误认为 model 本身能力。

\par \textbf{Multi-Agent System：}把一个任务交给多个 agents，各自拥有 role/context，并通过 communication/orchestration 协作：
\(
\text{Task}\rightarrow A_1,\ldots,A_K\rightarrow\text{Aggregator/Synthesis}.
\)
优势包括 parallel exploration、task specialization、独立 context windows；代价包括更多 tokens、communication overhead、duplicate work、conflicting answers 和 error propagation。
Multi-agent 的价值来自可以把相对独立的 subtasks 并行处理，然后 synthesis；如果任务高度 coupled、communication cost 很高或多个 agents 只是重复相同 reasoning，则
\(
\text{more agents}\not\Rightarrow\text{better performance}.
\)

\par \textbf{GPTSwarm：}把 LLM agent systems 表示为 \textbf{optimizable computational graphs}：nodes 表示 operations，edges 表示 information flow；不仅能优化 individual prompts，还可以通过改变 graph connectivity 自动优化 agent orchestration。
\textbf{GPTSwarm Hierarchy：}
Operation=Node,
Agent=Graph of Nodes,
Swarm=Composite Graph.
Agent 内 edges 定义 execution topology；agents 之间 edges 定义 collaboration/communication channels。

\par \textbf{GPTSwarm Graph Definition：}单个 agent 表示为 directed computational graph
\(
G=(N,E,F,o),
\)
其中 \(N\) 是 computational nodes，\(E\subseteq N\times N\) 是 directed edges，\(F=\{f_n\}_{n\in N}\) 是 node routines，\(o\) 是 output node；论文主要研究 DAG。Node 可以执行 LLM inference、tool use、function call 或 embodied action。
\textbf{GPTSwarm Node Execution：}node \(n\) 接收原始 input \(x\) 和 predecessor outputs \(z_n\)，执行
\(
f_n(z_n,x),
\)
并把结果传给 successors；最终 graph output 为
\(
\hat y=G(x)=f_o(z_o,x).
\)
因此整个 agent workflow 可以理解为一个显式 dataflow graph。

\par \textbf{GPTSwarm Composite Graph：}多个 agent graphs 通过 inter-agent edges 连接成 swarm；edge 不只是“两个 agent 可以聊天”，而是明确规定 information 从哪个 operation 流向哪个 operation，因此 communication topology 本身成为可以学习和优化的对象。

\par \textbf{GPTSwarm Node Optimization：}修改单个 node 的 prompt/behavior，使 local operation 更有效，即
\(
\text{node optimization}=\text{optimize what each component does}.
\)

\par \textbf{GPTSwarm Edge Optimization：}改变 graph connectivity，即决定哪些 nodes/agents 应该交换 information：
\text{edge optimization}=\text{optimize orchestration}.
论文的核心区别是，不只手工设计 multi-agent workflow，而是允许 communication structure 自动改进。
\textbf{GPTSwarm Edge Search：}若有 \(d\) 个 possible edges，则可能 graph 数量为
\(
2^d,
\)
无法 exhaustive search；因此定义 parameterized graph distribution
\(
G'\sim D_\theta
\)
并优化 expected task utility：
\(
\max_\theta
\mathbb E_{G'\sim D_\theta}
[u_\tau(G')].
\)

\par \textbf{REINFORCE for Graph Optimization：}因为 edge existence 和最终 task correctness 都是 discrete，不能直接普通 backprop；可以用 policy-gradient estimator：
\(
\nabla_\theta \mathbb E_{G\sim D_\theta}[u(G)]
=
\mathbb E[
u(G)\nabla_\theta\log p_\theta(G)
].
\)
高 reward graph 中出现的 communication structures 会被提高 probability，低 reward structures 则降低。

\par \textbf{GPTSwarm vs. Free-form Multi-Agent Chat：}普通 multi-agent system 常由人手工设计谁和谁聊天；GPTSwarm 把 communication 显式变成 graph edges，因此 workflow 可组合、可观察、可优化。CoT、Tree-of-Thought、Reflexion 等也可以看成不同 computational graph/topology。

\par \textbf{Function Calling：}开发者用 JSON Schema 描述可调用函数及参数，模型负责判断是否需要调用函数，并输出符合 function signature 的结构化 JSON arguments

\par \textbf{Function Calling vs. Tool Execution：}模型输出的 function call 只是 structured action request，并不会自己真正执行函数；完整流程是
\(
\text{LLM}
\rightarrow
\text{tool call}
\rightarrow
\text{Harness executes tool}
\rightarrow
\text{tool result}
\rightarrow
\text{LLM}.
\)
因此 “模型会 function calling” 和 “整个 agent 能可靠使用 tool” 不是一回事，后者还依赖 routing、permissions、error handling、retry 和 result injection。

\par \textbf{MCP (Model Context Protocol)：}Toolformer / LangChain-era 中，不同 framework 往往各自定义 tool description 和 calling protocol；MCP 将模型与 external tools / data sources 之间的接口标准化：
\(
\text{Model} \leftrightarrow \text{standard protocol} \leftrightarrow \text{Tools/Data Sources}.
\)
核心意义是让不同模型、client 和 tool servers 可以通过统一协议交互，而不需要每个 integration 都重新设计一套接口。

\par \textbf{Context-Window Tax：}暴露给模型的每一个 tool definition 都需要作为 context 的一部分输入，因此会消耗 context tokens：
\(
\text{more tools}
\Rightarrow
\text{more tool-schema tokens in context}.
\)
小而稳定的 tool set 成本较低；large tool catalog 不仅增加 token/context cost，也可能因为候选工具过多而降低 tool-selection accuracy。因此“让模型看到更多 tools”不是免费的。

\par \textbf{Stateful Session：}旧式 MCP design 中，每个 client 维持一个 live session，后续 request 依赖之前保存在某个 server instance 上的 session state；因此 load balancer 必须使用 sticky sessions，把同一 client 的后续 request 路由回原 server，或者所有 servers 使用 shared session store（如 Redis）。
\textbf{Stateless Request：}July 2026 MCP revision 改为每个 request 都是 self-contained，会携带自己的 protocol version、client identity 和 capabilities，因此任何 server instance 都可以独立处理任意 request，不依赖某个长期保存的 per-client session。
Stateful depends on server-side session state from previous requests
而
Stateless: each request contains enough information to be handled independently。
Stateless 的代价是需要在 request 中携带必要 metadata/context，但 infrastructure 更容易 horizontal scaling。

\par \textbf{Why Stateless Helps Load Balancing：}Stateful service 需要
\(
\text{client}\rightarrow\text{same server instance}
\)
的 sticky routing，或者额外维护 shared session store；stateless service 则可以让普通 round-robin load balancer 把任意 request 发送给任意 healthy server：
\(
R_1\rightarrow S_2,\quad
R_2\rightarrow S_5,\quad
R_3\rightarrow S_1.
\)
因此 server instances 更 interchangeable，扩容、缩容和故障恢复更简单。

\par \textbf{Why Stateless Helps Serverless：}Lambda、Cloudflare Workers 等 serverless infrastructure 不保证两次 requests 会落到同一个 persistent process，因此不适合依赖 local in-memory session state；stateless request 不依赖某个固定 server instance，所以可以自然地 scale up/down，并且不需要维护 shared session state。Slides 总结为：可以放在普通 round-robin load balancer 后面，也可以直接运行在 serverless infrastructure 上。

\par \textbf{Tool Schema Design：}tool schema 本身也是 harness design 的一部分；minimal schema 减少 context-window tax，但信息太少可能让模型不知道何时或如何调用；verbose schema 提供更多 guidance，却增加 context cost 和 ambiguity，因此需要在 \textbf{clarity vs. context cost} 之间权衡。

\topic{More}
Convolution 输出尺寸：
H\_out = floor((H + 2P - D(K-1) - 1)/S) + 1

Binary cross-entropy：
BCE = -[y log p + (1-y) log(1-p)]

\par \textbf{Why Not Single-Stream in ViLBERT：}一个直接方案是把 visual inputs discretize 成 visual tokens，与 text tokens 一起送进同一个 pretrained BERT；ViLBERT 指出三点问题：
(1) visual clustering / discretization 会引入 discretization error，丢失重要 visual details；
(2) image 与 language 被完全相同地处理，但两种 modality 的 complexity 和 input abstraction level 不同，因此可能需要不同 processing depth；
(3) 强迫 pretrained BERT weights 同时适应大量新的 visual tokens，可能破坏原本已经学好的 language representations。
因此 ViLBERT 使用 two-stream architecture，让 vision / language 分别处理，再通过少量 co-attention layers 交互。

\par \textbf{Another Single-Stream Cost：}single-stream 中 image 和 text 从很早开始一直交互，因此不能独立 cache 某一 modality 的 representations；ViLBERT 论文因此指出这种 architecture 在 image retrieval 等需要重复比较大量 image/text pairs 的任务上 computational cost 很高。

\par \textbf{Q-Former：}使用一组 learned query tokens，对 frozen vision encoder（如 CLIP / ViT）输出的 image features 做 cross-attention：
\(
Q=Q_{\mathrm{learned}},
K,V=\mathrm{VisionEncoder}(x).
\)
learned queries 主动从大量 visual features 中提取 task-relevant information，最后得到一个 compact image representation，作为 vision 与 language model 之间的 bridge。

\par \textbf{Why Q-Former：}vision encoder 往往输出很多 image patch / region features，而 LLM 不需要直接处理全部视觉 tokens；Q-Former 用少量 learned queries 抽取最有用的信息，因此既完成 visual-language alignment，又减少送入 LLM 的 token 数量。

\par \textbf{Q-Former Three Pretraining Objectives：}
(1) \textbf{Image-Text Matching (ITM)：}判断 image 和 text 是否对应；
(2) \textbf{Image-Text Contrastive Learning (ITC)：}让 matching image/text representations 靠近、non-matching pairs 远离；
(3) \textbf{Image-Grounded Text Generation：}根据 extracted image features condition text generation。
这三个 objectives 分别学习 matching、alignment 和 generation 能力。
\textbf{BLIP-2 Q-Former Bridge：}BLIP-2 先使用 frozen vision encoder 提取 image features，再通过 Q-Former 将 visual information 压缩成少量 representations，最后通过 projection / adapter ``glue'' 到 frozen LLM；这样不必 end-to-end retrain 巨大的 vision encoder 和 LLM。

\par \textbf{Q-Former vs Resampler：}二者都使用 learned queries + cross-attention 压缩 visual features，但重点不同：Q-Former 主要为 vision-language alignment 提取 task-relevant visual representation；Resampler 更强调把 variable-length visual embeddings 压缩成固定数量 tokens，以控制 LLM input size 和 computational cost。
\par \textbf{Resampler：}给定任意数量的 visual embeddings，通过固定数量 learned latent queries 做 cross-attention，输出 fixed-size token set；因此无论原图产生多少 patches / frames，送入 multimodal Transformer 的视觉 token 数量都保持固定。

\par \textbf{Early Fusion：}vision / text 在进入主要 Transformer self-attention 之前就已经被转换到统一 token / representation space，并在同一个 backbone 中 jointly processed；slides 中区分 Non-Tokenized Early Fusion (NTEF) 和 Tokenized Early Fusion (TEF)。
\textbf{Deep Fusion：}vision 与 language 先各自进行一定程度 processing，再在 Transformer 中间层通过 cross-attention 或 custom multimodal layers 反复融合；slides 将其分为 Standard Cross-attention based Deep Fusion (SCDF) 与 Custom Layer based Deep Fusion (CLDF)。
\textbf{Modern Fusion Trend：}slides 将 CLIP 归为 late fusion、ViLBERT 视为中间层 co-attention fusion；现代大型 multimodal LLM 更常采用 early fusion，把 image/video/text tokens 尽早放入同一 Transformer backbone，但 Q-Former / bridge-style deep fusion 仍常见于较便宜或 modular 的模型。

\par \textbf{LDM Text-to-Image Training：}给 training pair \((x,y)\)，其中 \(x\) 是 image、\(y\) 是 text prompt：
(1) encode image：
\(
z_0=E(x);
\)
(2) sample timestep \(t\)；
(3) sample Gaussian noise：
\(
\epsilon\sim\mathcal N(0,I);
\)
(4) 构造 noisy latent：
\(
z_t=\sqrt{\bar\alpha_t}z_0+
\sqrt{1-\bar\alpha_t}\epsilon;
\)
(5) encode text：
\(
c=\tau_\theta(y);
\)
(6) U-Net 根据 noisy latent、timestep 和 text condition 预测 noise：
\(
\hat\epsilon=\epsilon_\theta(z_t,t,c);
\)
(7) 使用 noise-prediction MSE：
\(
L=
\|\epsilon-\hat\epsilon\|_2^2;
\)
(8) gradient descent 更新 denoising model。一次 training iteration 只随机训练一个 timestep，不需要完整执行 \(z_T\rightarrow\cdots\rightarrow z_0\)。

\par \textbf{LDM Text-to-Image Generation：}给定 prompt \(y\)：
(1) encode text：
\(
c=\tau_\theta(y);
\)
(2) 从 pure Gaussian latent 开始：
\(
z_T\sim\mathcal N(0,I);
\)
(3) 对 \(t=T,\ldots,1\)，U-Net repeatedly predicts noise
\(
\hat\epsilon_t=\epsilon_\theta(z_t,t,c)
\)
并更新
\(
z_t\rightarrow z_{t-1};
\)
(4) 最终得到 clean latent \(z_0\)；
(5) decoder 只在最后执行一次：
\(
x=D(z_0).
\)
训练时从真实 image latent 加 noise；生成时没有原图，而是从 random noise 反向 denoise。

\par \textbf{DDPM：}Denoising Diffusion Probabilistic Model；forward process：
\(
x_t=
\sqrt{\bar\alpha_t}x_0+
\sqrt{1-\bar\alpha_t}\epsilon,
\qquad
\epsilon\sim\mathcal N(0,I).
\)
reverse process 是 stochastic：
\(
x_{t-1}
=
\mu_\theta(x_t,t)
+
\sigma_t\epsilon',
\qquad
\epsilon'\sim\mathcal N(0,I).
\)
因此即使从同一个 \(x_t\) 和 condition 开始，只要 reverse sampling noise 不同，最后生成的 image 也可能不同。

\par \textbf{Why DDPM Is Probabilistic：}reverse transition 建模为 distribution
\(
p_\theta(x_{t-1}\mid x_t),
\)
每一步不仅使用 model prediction，还额外 sample random Gaussian noise，因此 generation 是 stochastic trajectory。

\par \textbf{DDIM：}Denoising Diffusion Implicit Model；使用与 DDPM 相同或兼容的 trained noise predictor，但改变 sampling update，使 reverse process 可以减少甚至去掉额外随机 noise；一般形式可理解为
\(
x_{t-1}
=
\text{predicted clean-image component}
+
\text{predicted noise direction}
+
\sigma_t\epsilon.
\)
当 \(\sigma_t=0\) 时，sampling trajectory 可以变成 deterministic。

\par \textbf{DDPM vs DDIM：}DDPM reverse step 显式重新 sample noise，因此 stochastic；DDIM 可以令 reverse update deterministic，因此同一个 initial \(x_T\) + condition 会得到相同 output，并允许跳过部分 timesteps，从而通常可以用更少 sampling steps、生成更快。

\par \textbf{DDIM Inversion：}因为 deterministic DDIM 可以建立近似可逆 trajectory，可以把真实 image \(x_0\) 沿 deterministic path 映射到某个 noisy latent \(x_T\)，之后再从该 latent 反向 generation；这对 image editing / semantic manipulation 很有用，因为可以尽量保留原图 structure。

\par \textbf{Why DALL-E 2 Prior Is a Distribution：}给定同一个 caption \(y\)，可能存在很多不同但都合理的 images，因此 prior 不应该学 deterministic mapping
\(
y\rightarrow z_i,
\)
而应学习
\(
P(z_i\mid y),
\)
即从 text condition 产生一个 distribution over possible CLIP image embeddings；semantic diversity 很大程度来自这个 \(z_i\) distribution。

\par \textbf{Autoregressive Prior in DALL-E 2：}CLIP image embedding 本来是 continuous vector；autoregressive prior 先通过 PCA 降维，再将各 dimensions quantize 成 discrete tokens
\(
(c_1,c_2,\ldots,c_m),
\)
然后 Transformer autoregressively 建模
\(
P(c_1,\ldots,c_m\mid y)
=
\prod_jP(c_j\mid c_{<j},y).
\)

\par \textbf{Diffusion Prior in DALL-E 2：}直接把 CLIP image embedding \(z_i\) 当作 continuous data，在 embedding space 中加入 Gaussian noise：
\(
z_i\rightarrow z_i^{(1)}\rightarrow\cdots\rightarrow z_i^{(T)},
\)
再训练 Transformer-style diffusion model conditioned on text 来恢复 clean image embedding；本质是 ``diffusion in CLIP embedding space''，而不是 pixel diffusion。

\par \textbf{DALL-E 2 Two Sources of Diversity：}第一层 diversity 来自 prior：
\(
y\rightarrow z_i^{(1)},z_i^{(2)},\ldots
\)
同一个 prompt 可以得到不同 semantic interpretations；第二层 diversity 来自 decoder：
\(
z_i\rightarrow x^{(1)},x^{(2)},\ldots
\)
同一个 CLIP image embedding 又可以对应多张 pixel-level details 不同的 images。

\par \textbf{DALL-E 2 Decoder Resolution：}diffusion decoder 不是直接生成 \(1024\times1024\) image，而是先生成 \(64\times64\)，再用两个 diffusion upsamplers：
\(
64\times64
\rightarrow
256\times256
\rightarrow
1024\times1024.
\)
upsamplers 在已有低分辨率 semantic structure 基础上补充 high-frequency / spatial details，而不是重新决定主要 semantic content。

\par \textbf{CLIP Bottleneck in DALL-E 2：}DALL-E 2 generation 必须经过 CLIP image embedding \(z_i\)，因此 decoder 只能恢复 \(z_i\) 中保留下来的 information；如果 CLIP embedding 忽略某种细节，该信息在后续 decoder 中很难被准确恢复。

\par \textbf{Attribute Binding Limitation：}CLIP 更擅长表示整体 semantic similarity，但可能不能准确保存 ``哪个 attribute 属于哪个 object''，例如 ``red cube and blue sphere'' 可能混淆颜色与对象绑定，因此 DALL-E 2 也继承这种 attribute-binding weakness。

\par \textbf{Text / Spelling Limitation：}CLIP embedding 主要保存 high-level semantics，而不是精确 character-level spelling information，因此基于 CLIP latent 的 generation 对图片中的精确文字生成较弱。

\par \textbf{CLIP Compression Trade-off：}CLIP image embedding 是强 semantic bottleneck：优点是丢掉 irrelevant pixel details，使 image variation / semantic manipulation 更容易；缺点是被压掉的信息之后无法可靠重建。

% --- COLUMN 3 ---

\end{multicols}
\end{document}